% TOP_PROBLEM: {"id": 1193, "title": "Infinitude of Leinster Groups", "category_id": 4, "category": {"id": 4, "name": "algebra", "display_name": "Algebra", "description": "Group theory, ring theory, field theory, and algebraic structures.", "slug": "algebra", "order_index": 4, "created_at": "2026-07-31T15:26:25.671Z"}, "status": "open"}
% FIRST_POSTED: 2026-10-01
% ATTEMPT_STATUS: unresolved
% Imported from: unsolvedmath_additional_500.tex; UnsolvedMath ID 1193
% !TeX program = lualatex
% Compile twice: lualatex 1193.tex
% Self-contained: no external bibliography, images, or \input files.
% Numeric section labels and visible numbers are original UnsolvedMath IDs.
\documentclass[11pt,a4paper]{article}
\usepackage[margin=24mm,headheight=14pt]{geometry}
\usepackage{fontspec}
\setmainfont{Latin Modern Roman}
\setsansfont{Latin Modern Sans}
\setmonofont{Latin Modern Mono}
\usepackage{amsmath,amssymb,mathtools}
\usepackage{unicode-math}
\setmathfont{Latin Modern Math}
\usepackage{microtype}
\usepackage{enumitem,xurl,needspace}
\usepackage[dvipsnames]{xcolor}
\usepackage{hyperref}
\hypersetup{unicode=true,colorlinks=true,linkcolor=MidnightBlue,urlcolor=MidnightBlue,
 pdftitle={UnsolvedMath Problem 1193},pdfauthor={Source-annotated compilation},bookmarksnumbered=true}
\usepackage{bookmark,tocloft,titlesec,fancyhdr}
\setlength{\cftsecnumwidth}{4.2em}
\cftsetpnumwidth{2.5em}
\cftsetrmarg{4em}
\titleformat{\section}{\Large\bfseries\raggedright}{\thesection}{0.7em}{}
\pagestyle{fancy}\fancyhf{}
\fancyhead[L]{\small\sffamily Additional UnsolvedMath statements}
\fancyhead[R]{\small Original catalogue IDs}
\fancyfoot[C]{\thepage}
\setlength{\parindent}{0pt}
\setlength{\parskip}{5pt plus 1pt}
\setlength{\emergencystretch}{3em}
\setcounter{tocdepth}{1}\setcounter{secnumdepth}{1}
\setlist[itemize]{leftmargin=3em,itemsep=4pt,topsep=4pt,parsep=0pt}
\allowdisplaybreaks
\newcommand{\N}{\mathbb{N}}
\newcommand{\Z}{\mathbb{Z}}
\newcommand{\Q}{\mathbb{Q}}
\newcommand{\R}{\mathbb{R}}
\newcommand{\C}{\mathbb{C}}
\newcommand{\F}{\mathbb{F}}
\newcommand{\A}{\mathbb{A}}
\newcommand{\PP}{\mathbb{P}}
\newcommand{\E}{\mathbb{E}}
\newcommand{\eps}{\varepsilon}
\newcommand{\dd}{\,\mathrm d}
\newcommand{\sref}[2]{\hyperlink{source:#1:#2}{[S#2]}}
\newif\ifshowcatalogue
\showcataloguetrue
\newcommand{\cataloguescope}[1]{\ifshowcatalogue
 \paragraph{Statement recorded in the supplied source.}
 #1\fi}
\begin{document}
% UnsolvedMath ID 1193; source catalogue code ALG-029

\setcounter{section}{1192}

\section{Infinitely Many Leinster Groups}\label{1193}

{\small\sffamily UnsolvedMath ID 1193 · Algebra}

\subsection{Definitions and mathematical statement}

\paragraph{Shared notation.}Unless a section says otherwise, $\N=\{1,2,\ldots\}$,
$\Z$ is the ring of integers, $\Q,\R,\C$ are the rational, real, and complex
fields, $[N]=\{1,\ldots,\lfloor N\rfloor\}$, and $\log$ is the natural
logarithm. For real functions of a parameter tending to infinity,
$f\ll g$ and $f=O(g)$ mean $|f|\leq Cg$ eventually for some fixed $C$;
$f\gg g$ reverses this comparison; $f=o(g)$ means $f/g\to0$;
$f\sim g$ means $f/g\to1$; and $f\asymp g$ means both order bounds.
A subscript on $O$ or $\ll$ specifies allowed constant dependence.
The expression $n^{o(1)}$ in an upper bound means growth smaller than
$n^\epsilon$ for each fixed $\epsilon>0$. Local definitions govern any
exception, finite-field convention, or use of the same symbol for another
object. An asymptotic claim is not automatically an effective algorithm.



For a finite group $G$, let $\sigma_N(G)=\sum_{H\trianglelefteq G}|H|$, including both the identity subgroup and $G$. A Leinster group satisfies $\sigma_N(G)=2|G|$. Infinitude means infinitely many isomorphism classes of finite groups with this property, not infinitely many differently labelled copies of one group. \sref{1193}{1} \sref{1193}{2}

\paragraph{Statement recorded in the supplied source.}
Are there infinitely many Leinster groups? \sref{1193}{1}

\paragraph{Residual task reported by the archive.}

The embedded review (2026-08-17) records: Construct an infinite family or prove only finitely many exist. \sref{1193}{1}

\subsection{Short English statement}

Are there infinitely many finite groups for which the orders of all normal subgroups add up to twice the group’s order?

\subsection{Research attempt: coprime products and arithmetic obstructions}
\textbf{Status.} No infinite family is obtained. The inspected abstract of Leinster's original paper [S2] connects cyclic examples with perfect numbers. We prove that connection and a product criterion, then identify precisely why they do not yield unconditional infinitude.

\subsubsection{1. Cyclic groups and a conditional infinite supply}
A cyclic group $C_n$ has one subgroup of each divisor order $d\mid n$, all normal. Hence $\sigma_N(C_n)=\sigma(n)=\sum_{d\mid n}d$, and $C_n$ is Leinster exactly when $n$ is a perfect number. If $q=2^p-1$ is prime, then $p>1$, $q$ is odd, and coprime divisor factorization gives
\[\sigma(2^{p-1}q)=(1+2+\cdots+2^{p-1})(1+q)
 =(2^p-1)2^p=2(2^{p-1}q).\]
Thus infinitely many Mersenne primes would imply infinitely many pairwise nonisomorphic Leinster groups. The implication is proved; its antecedent is not proved here. The first orders from $p=2,3,5,7$ are $6,28,496,8128$, which the diagnostic verifies by direct divisor sums.

\subsubsection{2. Normal subgroup sums multiply under a precise hypothesis}
Let $G,H$ have coprime orders. Every subgroup $L\leq G\times H$ is a product. To see this, for $(g,h)\in L$ choose an integer $e$ congruent to 1 modulo $|G|$ and 0 modulo $|H|$. Then $(g,h)^e=(g,1)\in L$; a complementary exponent gives $(1,h)\in L$. Thus $L=L_G\times L_H$. Such a product is normal exactly when each factor is normal, as follows by conjugating separately in each coordinate. Consequently
\[\sigma_N(G\times H)=\sigma_N(G)\sigma_N(H),\qquad
\rho(G\times H)=\rho(G)\rho(H),\quad
\rho(G)=\frac{\sigma_N(G)}{|G|}.\tag{1}\]
This supplies an exact construction test: coprime factors must have ratios multiplying to 2. It does not allow repeated multiplication of existing examples; the product of two coprime Leinster groups has ratio 4.

\subsubsection{3. Excluding tempting factors}
No nontrivial finite $p$-group is Leinster. Its identity subgroup contributes 1 to $\sigma_N$, whereas every other subgroup order is divisible by $p$. Thus $\sigma_N\equiv1\pmod p$, while $2|G|\equiv0\pmod p$. For a nonabelian finite simple group $S$, only $1$ and $S$ are normal, so $\rho(S)=1+1/|S|<2$; it is not itself a solution either. Appending a coprime cyclic prime-order factor multiplies a ratio by $1+1/p>1$, and therefore cannot preserve ratio 2.

The coprimality hypothesis in (1) is indispensable. For example, $C_p\times C_p$ has $p+1$ subgroups of order $p$, so
\[\sigma_N(C_p\times C_p)=1+(p+1)p+p^2=1+p+2p^2,
\]
whereas $\sigma_N(C_p)^2=(1+p)^2$. Their difference is $p(p-1)>0$. The extra diagonal subgroups invalidate naive multiplicativity when orders overlap.

\subsubsection{4. Verification and what remains}
The diagnostic checks cyclic perfect orders below 10000, the displayed elementary-abelian calculation for $p=2,3,5,7$, and coprime cyclic multiplicativity for small orders. These tests confirm the arithmetic formulas rather than the existence of infinitely many groups. An unconditional infinite sequence of coprime ratio factorizations, or a construction controlling the additional normal subgroups when factors share primes, remains absent. Reducing infinitude to an unproved prime-producing assertion is not a solution of the original target.

\subsection{Sources}

{\small

\begin{itemize}

\item[\hypertarget{source:1193:1}{S1}] ULAM.AI, \emph{UnsolvedMath}, supplied \texttt{problems.json}, record 1193 (ALG-029), statement and background. The embedded literature-review date is 2026-08-17; that is the archive’s review date, not a fresh certification. Dataset landing page: \url{https://huggingface.co/datasets/ulamai/UnsolvedMath}.

\item[\hypertarget{source:1193:2}{S2}] T. Leinster, Perfect numbers and groups, arXiv:math/0104012 (2001). \url{https://arxiv.org/abs/math/0104012}. Bibliographic information imported from the supplied record.

\item[\hypertarget{source:1193:3}{S3}] T. Demedts et al., Perfect numbers and finite groups, preprint (). \url{https://cage.ugent.be/~tdemedts/preprints/leinster.pdf}. Bibliographic information imported from the supplied record.

\end{itemize}

}

\label{end:1193}
\end{document}
